% ============================================================================
% SPANISCH LANGENTWURF - SEQUENZ 5A: LOS HOBBYS
% ============================================================================
% Fachfarbe: #c41e3a (Karminrot)
% Tier 1 (vollständig)
% ============================================================================

\documentclass[11pt,a4paper]{article}

% ============================================================================
% PACKAGES
% ============================================================================
\usepackage[utf8]{inputenc}
\usepackage[T1]{fontenc}
\usepackage[ngerman]{babel}
\usepackage{geometry}
\usepackage{fancyhdr}
\usepackage{lastpage}
\usepackage{xcolor}
\usepackage{tcolorbox}
\usepackage{enumitem}
\usepackage{tabularx}
\usepackage{longtable}
\usepackage{array}
\usepackage{booktabs}
\usepackage{hyperref}
\usepackage{ifthen}
\usepackage{amssymb}

% ============================================================================
% KONFIGURATION
% ============================================================================

\newcommand{\plantier}{1}
\newcommand{\fachid}{spanisch}
\newcommand{\fach}{Spanisch}
\newcommand{\fachfarbe}{c41e3a}

% ============================================================================
% VARIABLEN
% ============================================================================

\newcommand{\jahrgangsstufe}{7}
\newcommand{\schulart}{Gesamtschule}
\newcommand{\stundenthema}{Los hobbys -- Hobbys benennen und erfragen}
\newcommand{\reihenthema}{Mi tiempo libre}
\newcommand{\stundennummer}{5a}
\newcommand{\reihenstunden}{4}
\newcommand{\unterrichtsdatum}{XX.XX.XXXX}
\newcommand{\unterrichtszeit}{XX:XX -- XX:XX}
\newcommand{\raum}{XXX}

\newcommand{\lehrkraft}{N.N.}
\newcommand{\ausbildungsstand}{Lehrkraft}

\newcommand{\lerngruppengroesse}{24}
\newcommand{\maennlich}{12}
\newcommand{\weiblich}{12}
\newcommand{\divers}{0}

\newcommand{\gerniveau}{A1}
\newcommand{\lehrwerk}{Qué pasa 1}
\newcommand{\lehrwerkseite}{S. 70-72}

% ============================================================================
% LAYOUT
% ============================================================================
\geometry{left=2.5cm, right=2.5cm, top=2.5cm, bottom=2.5cm}
\definecolor{fach}{HTML}{\fachfarbe}
\definecolor{gray}{HTML}{666666}
\definecolor{lightgray}{HTML}{f5f5f5}
\definecolor{successgreen}{HTML}{2a7a4b}
\definecolor{warningorange}{HTML}{b35c00}

\tcbuselibrary{skins,breakable}

\newtcolorbox{infobox}[1][]{
    colback=lightgray,
    colframe=fach,
    fonttitle=\bfseries,
    title=#1,
    breakable
}

\newtcolorbox{scorebox}{
    colback=white,
    colframe=fach,
    fonttitle=\bfseries,
    title=Didaktik-Score,
    breakable
}

\pagestyle{fancy}
\fancyhf{}
\fancyhead[L]{\small\textcolor{gray}{\fach{} -- Jg. \jahrgangsstufe{} -- \stundenthema}}
\fancyhead[R]{\small\textcolor{gray}{Langentwurf Tier 1}}
\fancyfoot[C]{\small\thepage{} / \pageref{LastPage}}
\renewcommand{\headrulewidth}{0.4pt}
\renewcommand{\footrulewidth}{0pt}

% ============================================================================
% DOKUMENT
% ============================================================================
\begin{document}

% ============================================================================
% DECKBLATT
% ============================================================================
\begin{titlepage}
    \centering
    \vspace*{2cm}

    {\large\textcolor{gray}{\schulart}}\\[2cm]

    {\Huge\textcolor{fach}{\textbf{Langentwurf}}}\\[0.5cm]
    {\large Tier 1 -- Vollständig}\\[2cm]

    {\LARGE\textbf{\stundenthema}}\\[0.5cm]
    {\large im Rahmen der Unterrichtsreihe}\\[0.3cm]
    {\Large\textit{\reihenthema}}\\[2cm]

    \begin{tabular}{rl}
        \textbf{Fach:} & \fach{} (\gerniveau) \\[0.3cm]
        \textbf{Klasse:} & \jahrgangsstufe \\[0.3cm]
        \textbf{Lehrwerk:} & \lehrwerk, \lehrwerkseite \\[0.3cm]
        \textbf{Datum:} & \underline{\hspace{3cm}} \\[0.3cm]
        \textbf{Zeit:} & \underline{\hspace{3cm}} (Doppelstunde) \\[0.3cm]
        \textbf{Raum:} & \underline{\hspace{3cm}} \\[0.3cm]
    \end{tabular}

    \vfill

    {\large Lehrkraft: \lehrkraft}\\[0.3cm]
    {\normalsize\textcolor{gray}{\ausbildungsstand}}\\[1cm]

\end{titlepage}

% ============================================================================
% INHALTSVERZEICHNIS
% ============================================================================
\tableofcontents
\newpage

% ============================================================================
% 1. BEDINGUNGSANALYSE
% ============================================================================
\section{Bedingungsanalyse}

\subsection{Lerngruppe}
\begin{tabular}{ll}
    Kursgröße: & \lerngruppengroesse{} SuS (\maennlich{} m / \weiblich{} w / \divers{} d) \\
    Jahrgangsstufe: & \jahrgangsstufe \\
    Sprachniveau: & \gerniveau{} (GER) \\
    Fremdsprache: & 2. Fremdsprache (Anfänger, 1. Lernjahr) \\
\end{tabular}

\subsection{Lernvoraussetzungen}

Die SuS...
\begin{itemize}
    \item beherrschen die Konjugation der Verben auf \textit{-ar} (Unidad 1)
    \item kennen die Verben \textit{ser} und \textit{tener}
    \item können sich vorstellen und nach dem Befinden fragen
    \item haben erste Erfahrungen mit Ausspracheregeln
\end{itemize}

\subsection{Differenzierung (intern)}

\textbf{WICHTIG:} Keine Sternchen auf Schülermaterial!

\begin{tabular}{|l|p{4cm}|p{4cm}|p{4cm}|}
\hline
& \textbf{[B] Basis} & \textbf{[S] Standard} & \textbf{[E] Erweiterung} \\
\hline
\textbf{Ziel} & Berufsreife & Mittlere Reife & Gymnasium \\
\hline
\textbf{Inhalt} & Vokabeln rezeptiv verstehen, Lückentext mit Hilfe & Vokabeln aktiv produzieren, Dialog mit Vorlage & Freie Dialoge, Zusatzvokabeln \\
\hline
\textbf{Hilfen} & Bildkarten, Wortliste & Teils Stichworte & Nur Impulse \\
\hline
\end{tabular}

\subsection{Besonderheiten}
\begin{itemize}
    \item \textbf{LRS:} 2 SuS (Nachteilsausgleich: Zeitzugabe, vergrößertes AB)
    \item \textbf{DaZ:} 1 SuS (Wortschatz deutsch-spanisch vorab besprechen)
    \item \textbf{Leistungsstark:} 4 SuS (Zusatzaufgaben, Expertenfunktion)
\end{itemize}

% ============================================================================
% 2. SACHANALYSE
% ============================================================================
\section{Sachanalyse}

\subsection{Sprachliche Strukturen}

\subsubsection{Neuer Wortschatz: Hobbys}
\begin{tabular}{ll}
\textbf{jugar al fútbol} & Fußball spielen \\
\textbf{escuchar música} & Musik hören \\
\textbf{ver la tele} & fernsehen \\
\textbf{leer} & lesen \\
\textbf{nadar} & schwimmen \\
\textbf{montar en bici} & Fahrrad fahren \\
\textbf{tocar la guitarra} & Gitarre spielen \\
\textbf{chatear} & chatten \\
\textbf{jugar a videojuegos} & Videospiele spielen \\
\end{tabular}

\subsubsection{Grammatik}
\begin{itemize}
    \item \textbf{Verben auf -er und -ir:} \textit{leer, ver} (hier nur Infinitiv)
    \item \textbf{Konstruktion jugar a + Sportart:} \textit{jugar al fútbol} (a + el = al)
    \item \textbf{Konstruktion tocar + Instrument:} \textit{tocar la guitarra}
\end{itemize}

\subsubsection{Phonetik}
\begin{itemize}
    \item \textbf{Aussprache j:} [x] wie deutsches ch in ,,ach'' (jugar, videojuegos)
    \item \textbf{Aussprache gu:} [g] vor e/i (guitarra)
    \item \textbf{Betonung:} Wörter auf Vokal betonen vorletzte Silbe (música $\rightarrow$ mú-si-ca)
\end{itemize}

\subsection{Kontrastive Betrachtung}

\begin{tabular}{|l|p{5cm}|p{5cm}|}
\hline
& \textbf{Deutsch} & \textbf{Spanisch} \\
\hline
Artikel bei Sport & Fußball spielen (ohne Artikel) & jugar \textbf{al} fútbol (mit Artikel) \\
\hline
Präposition & \textbf{--} Gitarre spielen & tocar \textbf{la} guitarra \\
\hline
Fernsehen & fernsehen (1 Wort) & ver la tele (3 Wörter) \\
\hline
\end{tabular}

\subsection{Kulturelle Aspekte}
\begin{itemize}
    \item \textbf{Fútbol:} Nationalsport in Spanien und Lateinamerika
    \item \textbf{Musikstile:} Reggaetón, Flamenco als typisch spanischsprachige Musik
    \item \textbf{Freizeitgestaltung:} Jugendliche verbringen viel Zeit mit Familie und Freunden draußen
\end{itemize}

% ============================================================================
% 3. DIDAKTISCHE ANALYSE
% ============================================================================
\section{Didaktische Analyse}

\subsection{Stellung in der Sequenz}
Dies ist die \textbf{\stundennummer. Doppelstunde} der Sequenz ``\reihenthema'' (\reihenstunden{} Doppelstunden gesamt).

\begin{tabular}{|c|l|l|}
\hline
\textbf{Block} & \textbf{Thema} & \textbf{Status} \\
\hline
5a & \textbf{Los hobbys} -- Hobbys benennen und erfragen & \textit{aktuell} \\
\hline
5b & Verbos en -er/-ir -- Konjugation & folgt \\
\hline
5c & Me gusta -- Vorlieben ausdrücken & folgt \\
\hline
5d & Sequenztest Mi tiempo libre & folgt \\
\hline
\end{tabular}

\subsection{Kompetenzziele (Rahmenplan MV)}

\begin{itemize}
    \item \textbf{Kommunikativ (Sprechen):} Über Freizeitaktivitäten sprechen, Fragen stellen und beantworten
    \item \textbf{Sprachlich (Wortschatz):} Hobbys und Freizeitaktivitäten benennen
    \item \textbf{Interkulturell:} Typische Freizeitaktivitäten spanischsprachiger Jugendlicher kennenlernen
    \item \textbf{Methodisch:} Vokabellerntechniken anwenden (Bild-Wort-Zuordnung)
\end{itemize}

\subsection{Legitimation}
\textbf{Rahmenplan Spanisch MV (Kl. 7-10):}
\begin{itemize}
    \item Kompetenzbereich: Kommunikative Kompetenz -- Sprechen
    \item Standard: ,,Die Schülerinnen und Schüler können einfache Gespräche über vertraute Themen aus dem Alltag führen.''
    \item Themenfeld: ,,yo y mi entorno'' -- persönliche Interessen und Freizeitgestaltung
\end{itemize}

% ============================================================================
% 4. STUNDENZIELE
% ============================================================================
\section{Stundenziele}

\subsection{Grobziel}
Die SuS erweitern ihre \textbf{kommunikative Kompetenz} im Bereich Sprechen, indem sie Hobbys benennen, nach Hobbys fragen und über ihre eigenen Freizeitaktivitäten sprechen.

\subsection{Feinziele (operationalisiert)}
\begin{tabular}{|c|p{8cm}|c|c|}
\hline
\textbf{Nr.} & \textbf{Die SuS können...} & \textbf{AFB} & \textbf{Diff.} \\
\hline
FZ1 & die 9 Hobby-Vokabeln korrekt aussprechen und den Bildern zuordnen & I & alle \\
\hline
FZ2 & die Frage ,,¿Qué te gusta hacer?'' verstehen und sinnvoll beantworten & II & alle \\
\hline
FZ3 & einen kurzen Dialog über Hobbys mit einem Partner führen & II & [S]/[E] \\
\hline
FZ4 & weitere Hobbys selbstständig recherchieren und benennen & III & [E] \\
\hline
\end{tabular}

% ============================================================================
% 5. METHODISCHE ANALYSE
% ============================================================================
\section{Methodische Analyse}

\subsection{Phasierung (Doppelstunde = 2 $\times$ 45 Min.)}
\begin{enumerate}
    \item \textbf{1. Stunde:}
    \begin{itemize}
        \item Einstieg: Bildimpuls (Jugendliche bei verschiedenen Aktivitäten)
        \item Erarbeitung: Vokabeln mit Bildkarten einführen
        \item Übung 1: Chorisches Sprechen, Nachsprechen
        \item Übung 2: Zuordnungsübung (Bild -- Wort)
    \end{itemize}
    \item \textbf{2. Stunde:}
    \begin{itemize}
        \item Wiederholung: Vokabelquiz (Bild zeigen $\rightarrow$ Schüler nennen)
        \item Anwendung: Dialoge ,,¿Qué te gusta hacer?''
        \item Festigung: Arbeitsblatt bearbeiten
        \item Sicherung: Ergebnisse besprechen
    \end{itemize}
\end{enumerate}

\subsection{Medien}
\begin{itemize}
    \item Tafel/Whiteboard
    \item Lehrwerk: \lehrwerk, \lehrwerkseite
    \item Arbeitsblatt: AB-05a.pdf
    \item Lernpfad: LP-05a.html (optional, digital)
    \item Bildkarten: 9 Hobbys (laminiert, A5)
    \item Audio: Track zu Unidad 4 (Hörverstehen Hobbys)
\end{itemize}

% ============================================================================
% 6. VERLAUFSPLANUNG
% ============================================================================
\section{Verlaufsplanung}

\textbf{1. Stunde (45 Min.)}

\begin{longtable}{|p{1cm}|p{1.5cm}|p{4cm}|p{3cm}|p{2cm}|p{2cm}|}
\hline
\textbf{Zeit} & \textbf{Phase} & \textbf{Lehrerhandeln} & \textbf{Schülerhandeln} & \textbf{SF} & \textbf{Medien} \\
\hline
\endhead

0--5 & Einstieg & Begrüßung auf Spanisch, Bildimpuls: ,,¿Qué hacen estos jóvenes?'' & Vermutungen äußern (auch auf Deutsch erlaubt) & Plenum & Tafel/Beamer \\
\hline

5--15 & Input & Vokabeln mit Bildkarten einführen, Bedeutung klären, Aussprache vorsprechen & Zuhören, nachsprechen & Plenum & Bildkarten \\
\hline

15--25 & Übung 1 & Chorisches Sprechen anleiten, auf Aussprache achten (j, gu) & Nachsprechen im Chor und einzeln & Plenum & -- \\
\hline

25--35 & Übung 2 & Zuordnungsübung anleiten: Bildkarten an Tafel, Wortkarten zuordnen lassen & Wortkarten den Bildern zuordnen & Plenum/EA & Bild-/Wortkarten \\
\hline

35--40 & Sicherung & Korrekte Zuordnung besprechen, Tafelanschrieb ergänzen & Ins Heft übertragen & Plenum & Tafel \\
\hline

40--45 & Über\-leitung & Ausblick 2. Stunde: ,,Wie frage ich nach Hobbys?'' & Aufräumen, Pausengong & Plenum & -- \\
\hline
\end{longtable}

\vspace{1em}

\textbf{2. Stunde (45 Min.)}

\begin{longtable}{|p{1cm}|p{1.5cm}|p{4cm}|p{3cm}|p{2cm}|p{2cm}|}
\hline
\textbf{Zeit} & \textbf{Phase} & \textbf{Lehrerhandeln} & \textbf{Schülerhandeln} & \textbf{SF} & \textbf{Medien} \\
\hline
\endhead

0--5 & Wieder\-holung & Vokabelquiz: Bild zeigen, SuS nennen Wort & Vokabeln nennen & Plenum & Bildkarten \\
\hline

5--15 & Input & Frage einführen: ,,¿Qué te gusta hacer?'' + Antwort: ,,Me gusta [Infinitiv]'' & Zuhören, nachsprechen & Plenum & Tafel \\
\hline

15--25 & Anwen\-dung & Partnerarbeit anleiten: Gegenseitig nach Hobbys fragen & Dialoge führen & PA & Dialogkarten \\
\hline

25--35 & Festigung & AB-05a ausgeben, unterstützen, differenziert helfen & AB bearbeiten & EA & AB-05a \\
\hline

35--42 & Sicherung & Lösungen besprechen, häufige Fehler thematisieren & Vergleichen, korrigieren & Plenum & Tafel/ML \\
\hline

42--45 & Abschluss & Zusammenfassung, HA erteilen & HA notieren & Plenum & -- \\
\hline
\end{longtable}

\textbf{Legende:} EA = Einzelarbeit, PA = Partnerarbeit, GA = Gruppenarbeit, SF = Sozialform

% ============================================================================
% 7. TAFELANSCHRIEB
% ============================================================================
\section{Tafelanschrieb}

\begin{center}
\fbox{\parbox{0.9\textwidth}{
\textbf{\large Los hobbys}

\vspace{0.5em}

\begin{tabular}{ll}
jugar al fútbol & Fußball spielen \\
escuchar música & Musik hören \\
ver la tele & fernsehen \\
leer & lesen \\
nadar & schwimmen \\
montar en bici & Fahrrad fahren \\
tocar la guitarra & Gitarre spielen \\
chatear & chatten \\
jugar a videojuegos & Videospiele spielen \\
\end{tabular}

\vspace{0.5em}

\textbf{Frage:} ¿Qué te gusta hacer?

\textbf{Antwort:} Me gusta + Infinitiv

\textit{Ejemplo:} Me gusta nadar.
}}
\end{center}

% ============================================================================
% 8. ERWARTETE SCHWIERIGKEITEN
% ============================================================================
\section{Erwartete Schwierigkeiten}

\begin{tabular}{|p{5cm}|p{8cm}|}
\hline
\textbf{Schwierigkeit} & \textbf{Reaktion} \\
\hline
Aussprache \textit{j} [x] (jugar, videojuegos) & Chorisches Sprechen, Einzelkorrektur, Vergleich mit dt. ,,ch'' in ,,ach'' \\
\hline
Verwechslung \textit{jugar al} / \textit{tocar la} & Kontrastive Übung: Sport vs. Instrument, visuelle Merkhilfe \\
\hline
\textit{ver la tele} (unregelmäßig) & Erklären: ,,sehen'' = \textit{ver}, nur Infinitiv in dieser Stunde \\
\hline
Zeitproblem (zu viele Vokabeln) & Puffer: Aufgabe 4 auf AB als HA, Reduzierung auf 6 Kernvokabeln \\
\hline
Heterogenität & Differenzierte AB-Aufgaben, Zusatzmaterial für [E] \\
\hline
\end{tabular}

% ============================================================================
% 9. HAUSAUFGABE
% ============================================================================
\section{Hausaufgabe}

\begin{itemize}
    \item Lehrwerk S. 71, Aufgabe 3 (Vokabelübung)
    \item Vokabeln lernen: \textit{jugar al fútbol} bis \textit{jugar a videojuegos}
    \item Optional: Lernpfad LP-05a.html bearbeiten (digital)
    \item [E] Zusatz: 3 weitere Hobbys recherchieren und aufschreiben
\end{itemize}

% ============================================================================
% 10. ANLAGEN
% ============================================================================
\section{Anlagen}

\begin{enumerate}
    \item Arbeitsblatt (AB-05a.pdf)
    \item Musterlösung (ML-05a.pdf)
    \item Lehrerhinweise (LH-05a.pdf)
    \item Lernpfad (LP-05a.html)
    \item Bildkarten Hobbys (9 Stück, A5)
    \item Dialogkarten für Partnerarbeit
\end{enumerate}

% ============================================================================
% 11. DIDAKTIK-SCORE
% ============================================================================
\newpage
\section{Didaktik-Score}

\begin{scorebox}
\textbf{Bewertung der didaktischen Qualität nach 10 Dimensionen}\\
\textbf{Ziel-Score: $\geq$ 9.0 (Premium-Qualität)}

\vspace{1em}
\begin{tabular}{|c|l|c|c|p{4.5cm}|}
\hline
\textbf{\#} & \textbf{Dimension} & \textbf{Gew.} & \textbf{Score} & \textbf{Notiz} \\
\hline
1 & Differenzierung & 15\% & 9/10 & [B] rezeptiv, [S] produktiv mit Hilfe, [E] frei + Zusatz \\
\hline
2 & Taxonomie (AFB) & 10\% & 9/10 & I: Zuordnung, II: Dialog führen, III: Recherche \\
\hline
3 & Lernziel-Alignment & 10\% & 10/10 & RP MV explizit adressiert (yo y mi entorno) \\
\hline
4 & Aktivierung & 10\% & 9/10 & Chorisch, PA-Dialoge, Zuordnung = konstruktiv \\
\hline
5 & Scaffolding & 10\% & 9/10 & Bildkarten, Dialogkarten, Wortlisten \\
\hline
6 & Feedback-Qualität & 10\% & 9/10 & Lösungsbesprechung, Aussprachekorrektur \\
\hline
7 & Sprachliche Klarheit & 10\% & 9/10 & Altersgerecht, max. 15 Wörter pro Satz \\
\hline
8 & Motivation \& Relevanz & 10\% & 10/10 & Alltagsbezug (Hobbys), Jugendthemen \\
\hline
9 & Inklusion (UDL) & 10\% & 9/10 & Bild + Text + Audio, LRS-Anpassung \\
\hline
10 & Praktikabilität & 5\% & 9/10 & 90 Min. realistisch, Material vorhanden \\
\hline
\multicolumn{3}{|r|}{\textbf{GESAMT:}} & \textbf{9.2/10} & \\
\hline
\end{tabular}

\vspace{1em}
\textbf{Bewertungsschwellen:}
\begin{itemize}
    \item[$\checkmark$] \textcolor{successgreen}{$\geq$ 9.0: \textbf{FREIGABE} (Premium-Qualität)}
    \item[$\Box$] \textcolor{warningorange}{8.0--8.9: Iteration empfohlen}
    \item[$\Box$] 6.0--7.9: Überarbeitung erforderlich
    \item[$\Box$] $<$ 6.0: Grundlegende Neukonzeption
\end{itemize}

\vspace{1em}
\textcolor{successgreen}{\textbf{Status: FREIGABE}}
\end{scorebox}

\end{document}
